\chapter{Badania i testy systemu}
\label{rozdzial5}

Celem tego rozdziału jest weryfikacja poprawności działania wytworzonego oprogramowania oraz ocena stopnia realizacji założonych wymagań. Przeprowadzono szereg testów funkcjonalnych oraz wydajnościowych.

\section{Scenariusze testowe}

W celu sprawdzenia kluczowych funkcjonalności systemu opracowano następujące scenariusze testowe, odzwierciedlające typowe ścieżki użytkownika:

\subsection{Scenariusz 1: Konfiguracja nowego gospodarstwa domowego}
\begin{enumerate}
    \item Użytkownik A rejestruje nowe konto.
    \item Użytkownik A tworzy nowy Dom o nazwie ,,Mieszkanie Studenckie''.
    \item Użytkownik A generuje kod zaproszenia i przekazuje go Użytkownikowi B.
    \item Użytkownik B rejestruje się i dołącza do domu używając otrzymanego kodu.
    \item \textbf{Oczekiwany rezultat:} Obaj użytkownicy widzą ten sam pulpit domu i mają dostęp do wspólnych zasobów.
\end{enumerate}

\subsection{Scenariusz 2: Dodawanie i rozliczanie rachunku z OCR}
\begin{enumerate}
    \item Użytkownik wybiera opcję dodania rachunku i robi zdjęcie paragonu.
    \item System przetwarza obraz i proponuje kwotę oraz nazwę sklepu.
    \item Użytkownik zatwierdza dane i wybiera podział kosztów ,,po równo''.
    \item Rachunek pojawia się na liście jako ,,nieopłacony''.
    \item \textbf{Oczekiwany rezultat:} Kwota zostaje poprawnie odczytana, a dług zostaje przypisany do pozostałych domowników.
\end{enumerate}

\section{Testy funkcjonalne backendu}

Testy jednostkowe i integracyjne warstwy serwerowej zostały zaimplementowane przy użyciu pakietu \texttt{testing} oraz biblioteki \texttt{testify}.

Przykładowy wynik uruchomienia testów dla serwisu autoryzacji:
\begin{lstlisting}[language=bash, caption=Wynik testów jednostkowych, label=test:output]
$ go test ./internal/services/... -v
=== RUN   TestRegisterUser
--- PASS: TestRegisterUser (0.00s)
=== RUN   TestLoginUser_Success
--- PASS: TestLoginUser_Success (0.00s)
=== RUN   TestLoginUser_InvalidPassword
--- PASS: TestLoginUser_InvalidPassword (0.00s)
PASS
ok      github.com/Dragodui/diploma-server/internal/services    0.123s
\end{lstlisting}

Wszystkie kluczowe ścieżki krytyczne (rejestracja, tworzenie domu, dodawanie zadań) zostały pokryte testami automatycznymi, co zapewnia stabilność API.

\section{Testy interfejsu mobilnego}

Testy aplikacji mobilnej przeprowadzono manualnie na dwóch urządzeniach fizycznych:
\begin{itemize}
    \item Google Pixel 6 (Android 14)
    \item iPhone 13 (iOS 17)
\end{itemize}
oraz na emulatorach.

Weryfikowano:
\begin{itemize}
    \item \textbf{Responsywność:} Czy elementy UI skalują się poprawnie na różnych rozdzielczościach.
    \item \textbf{Nawigację:} Czy przechodzenie między ekranami jest płynne i intuicyjne.
    \item \textbf{Obsługę błędów:} Czy aplikacja wyświetla czytelne komunikaty w przypadku braku internetu lub błędu serwera.
\end{itemize}

Stwierdzono, że aplikacja działa stabilnie na obu platformach. Wykryte drobne błędy wizualne (np. przesunięcia tekstu) zostały na bieżąco poprawione.

\section{Wyniki i wnioski z testów}

Przeprowadzone badania wykazały, że system spełnia założone wymagania funkcjonalne.

\begin{itemize}
    \item \textbf{Wydajność:} Czas odpowiedzi API dla standardowych operacji wynosi średnio poniżej 100ms (w środowisku lokalnym).
    \item \textbf{Skuteczność OCR:} Moduł rozpoznawania tekstu poprawnie odczytuje kwoty z paragonów w około 80% przypadków przy dobrym oświetleniu. W przypadku słabej jakości zdjęcia, użytkownik ma możliwość ręcznej edycji danych.
    \item \textbf{Użyteczność:} Testy z udziałem użytkowników potwierdziły, że interfejs jest przejrzysty, a proces zakładania domu i zapraszania domowników nie sprawia trudności.
\end{itemize}

System jest gotowy do wdrożenia w wersji produkcyjnej (MVP).
